\begin{tikzpicture}[
  node distance=1.2cm,
  code/.style={
    draw,
    rounded corners,
    fill=gray!8,
    font=\ttfamily\small,
    align=left,
    inner sep=8pt
  },
  note/.style={
    draw,
    rounded corners,
    fill=blue!6,
    font=\small,
    align=left,
    text width=4.2cm,
    inner sep=6pt
  },
  arrow/.style={->, thick}
]

\node[code] (code) {
let D : ps.carrier $\to$ Prop :=\\
\quad fun candidate\_element =>\\
\quad\quad candidate\_element = ps.one $\lor$\\
\quad\quad $\exists$ predecessor : ps.carrier,\\
\quad\quad\quad ps.successor predecessor = candidate\_element
};

\node[note, right=2.8cm of code.north east, anchor=north west] (subset) {
\textbf{D is a subset predicate}\\
In Lean, a subset of the carrier is represented as a function:
\[
D : ps.carrier \to Prop.
\]
It answers: ``does this element belong to \(D\)?''
};

\node[note, right=2.8cm of code.east, anchor=west] (candidate) {
\textbf{candidate\_element}\\
This is the element currently being tested for membership in \(D\).
It plays the role of \(x\) in:
\[
D=\{x\in P : \cdots\}.
\]
};

\node[note, right=2.8cm of code.south east, anchor=south west] (condition) {
\textbf{Membership condition}\\
The proposition says:
\[
x = 1
\quad\text{or}\quad
\exists u,\ S(u)=x.
\]
So \(D\) contains exactly elements that are either \(1\) or successors.
};

\node[note, below=1.5cm of code] (predecessor) {
\textbf{predecessor}\\
The existential witness:
\[
\exists\, predecessor : ps.carrier
\]
means ``there is some element of the carrier whose successor is the candidate element.''
};

\draw[arrow] (subset.west) -- ++(-0.9,0) |- ([yshift=0.45cm]code.east);
\draw[arrow] (candidate.west) -- ++(-0.9,0) |- ([yshift=-0.15cm]code.east);
\draw[arrow] (condition.west) -- ++(-0.9,0) |- ([yshift=-0.95cm]code.east);
\draw[arrow] (predecessor.north) -- ([yshift=-1.2cm]code.south);

\end{tikzpicture}
